\begin{frame}[plain,fragile]
  \slideheader{Motivation}

  \vspace{15mm}
  {\small\textbf{Problem: classifier performance can degrade under domain shift.}}\par

  \vspace{2mm}
  \begin{itemize}\tightlist
    \item GN3/SALT is trained for jet-flavour classification.
    \item Inputs may come from multiple domains: nominal simulation/reconstruction, alternate samples, systematic variations, or shifted distributions.
    \item A model trained only to classify nominal data can learn representations that are not stable across domains.
    \item Goal: preserve flavour discrimination while making the learned embedding less sensitive to nominal--alternate differences.
  \end{itemize}

  \takeaway{Robust tagging needs representations that remain discriminative for flavour but less domain-specific.}

  \slideheaderfooter
\end{frame}

\begin{frame}[plain,fragile]
  \slideheader{Design Goal}

  \vspace{15mm}
  {\small\textbf{Minimum viable domain-alignment implementation.}}\par

  \vspace{2mm}
  \begin{columns}[T,totalwidth=\textwidth]
    \begin{column}{0.50\textwidth}
      \begin{itemize}\tightlist
        \item Use one shared GN3/SALT backbone.
        \item Keep the standard jet-classification head.
        \item Add Sinkhorn alignment on \texttt{global\_rep}.
        \item Treat alternate-domain samples as unlabeled for optimization.
      \end{itemize}
    \end{column}
    \begin{column}{0.46\textwidth}
      \begin{itemize}\tightlist
        \item No MACK/ROCK integration yet.
        \item No adversarial/SAT-style training yet.
        \item No auxiliary GN2/GN3 tasks in the MVP.
        \item No dynamic GLS/DWA loss weighting for the first study.
      \end{itemize}
    \end{column}
  \end{columns}

  \takeaway{This is a small, controlled SALT extension rather than a full architecture redesign.}

  \slideheaderfooter
\end{frame}

\begin{frame}[plain,fragile]
  \slideheader{Dataset Setup}

  \vspace{14mm}
  \begin{columns}[T,totalwidth=\textwidth]
    \begin{column}{0.45\textwidth}
      {\small\textbf{Existing mixed-domain H5 layout}}\par
      \begin{lstlisting}[style=textcode]
jets.source_label = 0 or 1
jets.flavour_label = class label
      \end{lstlisting}
      \vspace{1mm}
      \begin{lstlisting}[style=yamlcode]
classification_nominal_value: 0
nominal_value: 0
alternate_value: 1
      \end{lstlisting}
    \end{column}
    \begin{column}{0.51\textwidth}
      \begin{itemize}\tightlist
        \item \texttt{source\_label} identifies nominal vs alternate domain.
        \item \texttt{flavour\_label} is used only for nominal supervised classification.
        \item Alternate labels may exist in H5, but are ignored during optimization.
        \item \texttt{source\_label} is kept out of \texttt{data.variables} to avoid direct domain leakage.
      \end{itemize}
      \vspace{2mm}
      {\scriptsize\color{argonneRed}\textbf{Open check:} confirm semantic mapping of \texttt{source\_label == 0} vs \texttt{1} before production training.}
    \end{column}
  \end{columns}

  \takeaway{Domain labels are metadata for masking/splitting losses, not model inputs.}

  \slideheaderfooter
\end{frame}

\begin{frame}[plain,fragile]
  \slideheader{Training Objective}

  \vspace{15mm}
  \begin{center}
    {\large\textbf{$L_{\mathrm{total}} = L_{\mathrm{CE, nominal}} + \lambda_{\mathrm{Sinkhorn}}(t)\,D_{\mathrm{Sinkhorn}}(z_{\mathrm{nominal}}, z_{\mathrm{alternate}})$}}
  \end{center}

  \vspace{2mm}
  \begin{itemize}\tightlist
    \item \textbf{Nominal classification loss}: standard jet-classification CE, evaluated only on nominal rows.
    \item \textbf{Sinkhorn alignment loss}: compares nominal and alternate \texttt{global\_rep} embedding distributions.
    \item \textbf{Alignment strength}: configured by the Sinkhorn loss \texttt{weight}, with optional warmup/ramp scaling.
    \item \textbf{Diagnostics}: raw Sinkhorn and ramp scale are logged separately and excluded from the optimized total loss.
  \end{itemize}

  \takeaway{Trusted nominal labels drive classification; both domains shape a more stable representation.}

  \slideheaderfooter
\end{frame}

\begin{frame}[plain,fragile]
  \slideheader{Why Sinkhorn?}

  \vspace{15mm}
  {\small\textbf{Optimal transport gives a distribution-level alignment objective.}}\par

  \vspace{2mm}
  \begin{columns}[T,totalwidth=\textwidth]
    \begin{column}{0.52\textwidth}
      \begin{itemize}\tightlist
        \item Compares embedding point clouds directly.
        \item Does not require alternate-domain labels.
        \item Provides a smooth, differentiable objective.
        \item Works as a headless auxiliary loss on \texttt{global\_rep}.
      \end{itemize}
    \end{column}
    \begin{column}{0.44\textwidth}
      \begin{itemize}\tightlist
        \item No domain discriminator.
        \item No adversarial min--max loop.
        \item No second output head.
        \item Directly fits into the existing SALT training step.
      \end{itemize}
    \end{column}
  \end{columns}

  \takeaway{Sinkhorn alignment encourages domain-stable embeddings without adversarial machinery.}

  \slideheaderfooter
\end{frame}

\begin{frame}[plain,fragile]
  \slideheader{Implementation Overview}

  \vspace{14mm}
  \begin{columns}[T,totalwidth=\textwidth]
    \begin{column}{0.47\textwidth}
      {\small\textbf{Main additions inside \texttt{ModelWrapper}}}\par
      \begin{itemize}\tightlist
        \item Optional \texttt{domain\_alignment} module.
        \item Domain-label config for CE masking and embedding split.
        \item Optional nominal-only classification loss.
        \item Epoch-based warmup/ramp controls.
        \item Diagnostic logging for raw Sinkhorn and scale.
        \item Guard: initial Sinkhorn work uses \texttt{loss\_mode: wsum}.
      \end{itemize}
    \end{column}
    \begin{column}{0.49\textwidth}
      \begin{lstlisting}[style=textcode]
salt/modelwrapper.py
salt/models/SinkhornDomainAlignmentLoss.py
      \end{lstlisting}
      \vspace{1mm}
      \begin{lstlisting}[style=yamlcode]
domain_label: source_label
nominal_value: 0
alternate_value: 1
mask_alternate_classification_loss: true
      \end{lstlisting}
    \end{column}
  \end{columns}

  \takeaway{The Sinkhorn loss is integrated as an optional extension to the wrapper.}

  \slideheaderfooter
\end{frame}

\begin{frame}[plain,fragile]
  \slideheader{Model/Data Flow}

  \vspace{10mm}
  \centering
  \begin{tikzpicture}[x=1mm,y=1mm,
      box/.style={draw=argonneBlue, thick, rounded corners=1.5mm, fill=argonneBlue!6,
                  align=center, inner sep=1.4mm, minimum height=8mm},
      lossbox/.style={draw=argonneGreen!70!black, thick, rounded corners=1.5mm,
                      fill=argonneGreen!10, align=center, inner sep=1.4mm, minimum height=8mm},
      note/.style={align=center, font=\scriptsize},
      arrow/.style={-Latex, thick, argonneRed}]

    \node[box, minimum width=34mm] (batch) at (78,-2) {\textbf{Mixed H5 batch}\\[-0.3mm]\scriptsize \texttt{source\_label} + features};
    \node[box, minimum width=44mm] (backbone) at (78,-18) {\textbf{GN3/SALT backbone}\\[-0.3mm]\scriptsize features exclude \texttt{source\_label}};
    \node[box, minimum width=34mm] (rep) at (34,-34) {\textbf{global\_rep}\\[-0.3mm]\scriptsize shared embedding};
    \node[box, minimum width=40mm] (head) at (112,-34) {\textbf{classification head}\\[-0.3mm]\scriptsize jet flavour logits};
    \node[lossbox, minimum width=42mm] (sinkhorn) at (34,-50) {\textbf{Sinkhorn alignment}\\[-0.3mm]\scriptsize nominal vs alternate};
    \node[lossbox, minimum width=45mm] (ce) at (112,-50) {\textbf{CE loss}\\[-0.3mm]\scriptsize nominal rows only};

    \draw[arrow] (batch) -- (backbone);
    \draw[arrow] (backbone) -- (rep);
    \draw[arrow] (backbone) -- (head);
    \draw[arrow] (rep) -- (sinkhorn);
    \draw[arrow] (head) -- (ce);
    \node[note, text width=78mm, anchor=north] at (78,-58) {\texttt{source\_label} is used only after the forward pass for splitting embeddings and masking CE.};
  \end{tikzpicture}

  \slideheaderfooter
\end{frame}

\begin{frame}[plain,fragile]
  \slideheader{Warmup and Ramp}

  \vspace{11mm}
  \begin{columns}[T,totalwidth=\textwidth]
    \begin{column}{0.42\textwidth}
      {\small\textbf{Preserve a classification-only baseline.}}\par
      \begin{lstlisting}[style=yamlcode]
alignment_warmup_epochs: 100
alignment_ramp_epochs: 20
      \end{lstlisting}
    \end{column}
    \begin{column}{0.54\textwidth}
      \begin{itemize}\tightlist
        \item Epochs \texttt{0--99}: classification-only.
        \item Epoch \texttt{099}: warmup-end baseline checkpoint.
        \item Epoch \texttt{100+}: Sinkhorn alignment begins.
        \item Ramp gradually increases the alignment contribution.
      \end{itemize}
    \end{column}
  \end{columns}

  \vspace{2mm}
  \begin{center}
    \begin{tikzpicture}[x=1mm,y=1mm]
      \draw[->, thick] (0,0) -- (98,0) node[anchor=west]{\scriptsize epoch};
      \draw[->, thick] (0,0) -- (0,18) node[anchor=south]{\scriptsize scale};
      \draw[argonneBlue, very thick] (0,0) -- (50,0) -- (68,14) -- (94,14);
      \draw[dashed] (50,0) -- (50,16);
      \draw[dashed] (68,0) -- (68,16);
      \node[anchor=north] at (50,-1) {\scriptsize 100};
      \node[anchor=north] at (68,-1) {\scriptsize 120};
      \node[anchor=south] at (25,3) {\scriptsize warmup};
      \node[anchor=south] at (59,7) {\scriptsize ramp};
      \node[anchor=south] at (82,14) {\scriptsize full weight};
    \end{tikzpicture}
  \end{center}

  \takeaway{One training run can provide both baseline and domain-aligned checkpoints.}

  \slideheaderfooter
\end{frame}

\begin{frame}[plain,fragile]
  \slideheader{Configuration}

  \vspace{12mm}
  {\small\textbf{Main config:} \texttt{salt/configs/GN3/GN3\_sinkhorn\_domain\_alignment.yaml}}\par

  \begin{lstlisting}[style=yamlcode]
model:
  loss_mode: wsum
  mask_alternate_classification_loss: true
  classification_domain_label: source_label
  classification_nominal_value: 0
  alignment_warmup_epochs: 100
  alignment_ramp_epochs: 20
  domain_alignment:
    class_path: salt.models.SinkhornDomainAlignmentLoss
    init_args:
      input_name: jets
      domain_label: source_label
      nominal_value: 0
      alternate_value: 1
      weight: 0.1
      blur: 0.1
      scaling: 0.8
      normalization: layernorm
      min_samples_per_domain: 8
  \end{lstlisting}

  \takeaway{One config can run warmup, active Sinkhorn training, and checkpoint-based comparison.}

  \slideheaderfooter
\end{frame}

\begin{frame}[plain,fragile]
  \slideheader{Batch Composition Checks}

  \vspace{14mm}
  {\small\textbf{Sinkhorn requires both domains in each batch.}}\par

  \vspace{1mm}
  \begin{columns}[T,totalwidth=\textwidth]
    \begin{column}{0.52\textwidth}
      \begin{lstlisting}[style=bashcode]
python salt/utils/check_domain_batches.py \
  --file /path/to/pp_output_train.h5 \
  --dataset jets \
  --domain-label source_label \
  --batch-size 256 \
  --min-samples-per-domain 8 \
  --ignore-tail
      \end{lstlisting}
    \end{column}
    \begin{column}{0.44\textwidth}
      \begin{itemize}\tightlist
        \item Scans contiguous H5 windows matching SALT batch size.
        \item Reports min/max nominal and alternate counts.
        \item Flags single-domain batches.
        \item Flags unexpected domain-label values.
      \end{itemize}
    \end{column}
  \end{columns}

  \vspace{2mm}
  {\small\textbf{Current validation summary:} train and validation batches contain both domains; no single-domain batches or unexpected labels found.}\par

  \takeaway{The existing mixed H5 layout is compatible with per-batch Sinkhorn alignment.}

  \slideheaderfooter
\end{frame}

\begin{frame}[plain,fragile]
  \slideheader{Environment and Dependencies}

  \vspace{14mm}
  \begin{columns}[T,totalwidth=\textwidth]
    \begin{column}{0.45\textwidth}
      {\small\textbf{Lightweight repo-local environment}}\par
      \begin{lstlisting}[style=bashcode]
uv venv .venv-salt
source .venv-salt/bin/activate
uv sync
      \end{lstlisting}
      \vspace{1mm}
      \begin{lstlisting}[style=textcode]
geomloss>=0.2.6
      \end{lstlisting}
    \end{column}
    \begin{column}{0.51\textwidth}
      {\small\textbf{Smoke validation completed}}\par
      \begin{itemize}\tightlist
        \item PyTorch CUDA available.
        \item \texttt{geomloss} imports successfully.
        \item \texttt{salt --help} works.
        \item Short GPU training smoke test completes.
      \end{itemize}
    \end{column}
  \end{columns}

  \takeaway{The implementation adds one main dependency and works in a lightweight local environment.}

  \slideheaderfooter
\end{frame}

\begin{frame}[plain,fragile]
  \slideheader{Smoke Test Results}

  \vspace{14mm}
  {\small\textbf{Initial GPU smoke test passed.}}\par

  \vspace{1mm}
  \begin{columns}[T,totalwidth=\textwidth]
    \begin{column}{0.45\textwidth}
      \begin{itemize}\tightlist
        \item Config loads successfully.
        \item Data module initializes.
        \item Model includes \texttt{SinkhornDomainAlignmentLoss}.
        \item Training runs on GPU.
        \item Short 1-epoch smoke completed.
      \end{itemize}
    \end{column}
    \begin{column}{0.51\textwidth}
      {\small\textbf{Caveat and next smoke}}\par
      \begin{itemize}\tightlist
        \item Default warmup means a 1-epoch run mostly tests the classification path.
        \item FlashAttention missing warning appears, but SALT falls back to \texttt{torch-math}.
      \end{itemize}
      \begin{lstlisting}[style=bashcode]
salt fit -c salt/configs/GN3/GN3_sinkhorn_domain_alignment.yaml \
  --model.alignment_warmup_epochs 0 \
  --model.alignment_ramp_epochs 1
      \end{lstlisting}
    \end{column}
  \end{columns}

  \takeaway{The plumbing works; active-Sinkhorn smoke and full training are the next validation steps.}

  \slideheaderfooter
\end{frame}

\begin{frame}[plain,fragile]
  \slideheader{Checkpointing and Inference}

  \vspace{14mm}
  {\small\textbf{Per-epoch checkpoints are available.}}\par

  \vspace{1mm}
  \begin{columns}[T,totalwidth=\textwidth]
    \begin{column}{0.44\textwidth}
      \begin{itemize}\tightlist
        \item SALT checkpoint callback uses \texttt{save\_top\_k = -1}.
        \item Checkpoints are kept under the run directory.
        \item Warmup-end checkpoint is the classification-only baseline.
        \item Later checkpoints include active alignment.
      \end{itemize}
    \end{column}
    \begin{column}{0.52\textwidth}
      \begin{lstlisting}[style=textcode]
logs/<run_name>_<timestamp>/ckpts/
      \end{lstlisting}
      \begin{lstlisting}[style=bashcode]
salt test \
  -c logs/<run>/config.yaml \
  --ckpt_path logs/<run>/ckpts/epoch=099-val_loss=....ckpt
      \end{lstlisting}
      {\scriptsize If \texttt{--ckpt\_path} is omitted, SALT selects the lowest-validation-loss checkpoint from \texttt{ckpts/}.}\par
    \end{column}
  \end{columns}

  \takeaway{Baseline and aligned models can be compared from checkpoints produced by the same training run.}

  \slideheaderfooter
\end{frame}

\begin{frame}[plain,fragile]
  \slideheader{What This Enables}

  \vspace{14mm}
  \begin{columns}[T,totalwidth=\textwidth]
    \begin{column}{0.43\textwidth}
      {\small\textbf{Direct comparison}}\par
      \begin{enumerate}\tightlist
        \item Warmup-end checkpoint: classification-only baseline.
        \item Later checkpoints: same architecture and data, plus representation alignment.
      \end{enumerate}
    \end{column}
    \begin{column}{0.53\textwidth}
      {\small\textbf{Evaluation questions}}\par
      \begin{itemize}\tightlist
        \item Does nominal classification performance remain stable?
        \item Does alternate-domain performance improve?
        \item Do nominal and alternate \texttt{global\_rep} distributions become closer?
        \item How sensitive are results to Sinkhorn weight?
      \end{itemize}
    \end{column}
  \end{columns}

  \takeaway{The implementation creates a controlled setup for measuring domain-alignment benefit.}

  \slideheaderfooter
\end{frame}

\begin{frame}[plain,fragile]
  \slideheader{Planned Hyperparameter Scan}

  \vspace{14mm}
  {\small\textbf{Start with Sinkhorn strength.}}\par

  \vspace{2mm}
  \begin{columns}[T,totalwidth=\textwidth]
    \begin{column}{0.42\textwidth}
      \begin{lstlisting}[style=textcode]
weight = 0.0, 0.01, 0.05, 0.1
      \end{lstlisting}
      \begin{itemize}\tightlist
        \item \texttt{0.0}: no-alignment control.
        \item \texttt{0.01}: weak alignment.
        \item \texttt{0.05}: nominal starting point.
        \item \texttt{0.1}: stronger alignment.
      \end{itemize}
    \end{column}
    \begin{column}{0.54\textwidth}
      {\small\textbf{Monitor}}\par
      \begin{itemize}\tightlist
        \item nominal validation classification loss;
        \item alternate-domain test performance;
        \item flavour-tagging metrics;
        \item \texttt{sinkhorn\_loss} trends;
        \item embedding distribution diagnostics.
      \end{itemize}
    \end{column}
  \end{columns}

  \takeaway{Keep the first study simple: scan the primary alignment knob.}

  \slideheaderfooter
\end{frame}

\begin{frame}[plain,fragile]
  \slideheader{Current Limitations}

  \vspace{15mm}
  {\small\textbf{Known caveats before physics interpretation.}}\par

  \vspace{2mm}
  \begin{itemize}\tightlist
    \item Semantic mapping of \texttt{source\_label == 0} vs \texttt{source\_label == 1} still needs confirmation.
    \item Full production training has not yet been run.
    \item Scientific evaluation is still pending.
    \item FlashAttention is not installed, so training may be slower or more memory-intensive.
    \item No MACK/ROCK projector yet.
    \item No adversarial training or dynamic loss weighting included.
  \end{itemize}

  \takeaway{The implementation is ready for controlled studies, but physics validation is still required.}

  \slideheaderfooter
\end{frame}

\begin{frame}[plain,fragile]
  \slideheader{Summary}

  \vspace{15mm}
  {\small\textbf{Minimal Sinkhorn domain alignment is now integrated into SALT/GN3.}}\par

  \vspace{2mm}
  \begin{columns}[T,totalwidth=\textwidth]
    \begin{column}{0.50\textwidth}
      \begin{itemize}\tightlist
        \item Added optional Sinkhorn alignment on \texttt{global\_rep}.
        \item Masked classifier loss to nominal-domain rows.
        \item Kept alternate-domain rows unlabeled for optimization.
        \item Preserved standard GN3/SALT architecture.
      \end{itemize}
    \end{column}
    \begin{column}{0.46\textwidth}
      \begin{itemize}\tightlist
        \item Added one Sinkhorn-capable GN3 config.
        \item Added batch-composition checking utility.
        \item Verified basic GPU smoke behavior.
        \item Established a checkpoint-based baseline/aligned comparison strategy.
      \end{itemize}
    \end{column}
  \end{columns}

  \takeaway{This is a clean baseline for testing whether optimal-transport alignment improves domain robustness.}

  \slideheaderfooter
\end{frame}
